\section{Gestures}
\subsection{Introduction}
Hand gestures are an integral additional component when considering interpersonal communication. They allow for further nuances and meaning to be added to conventional communication~\citep{Sato01}. In addition, they are often able to provide a multitude of symbolic meaning in a very simple and brief manner. Due to their permeation of traditional communication, they are also an intuitive means of providing commands as their purpose and nature is familiar~\citep{Kett00}. With the advent of Computer Vision and touch-based systems, along with the increased penetration of computing devices in everyday life, gestures have begun to be explored as a means for interacting with computers~\citep{Rubi91}. The following review will discuss hand gestures in this context by exploring their benefits and limitations.

\subsection{Gestures versus direct manipulation}
The definition of gestures is often highly dependant on their application. However, an overarching theme is that they are symbolic movements used to relay scope and commands~\citep{Rubi91}. According to \citet{Huan95}, ``Every gesture is the physical expression of a mental concept". In the context of Human-Computer Interaction, gestures generally serve two purposes: they are used to draw attention to some point or area and to manipulate or act upon an object. These two functions are often linked within the same gesture, highlighting the context-sensitive nature of most gestural commands~\citep{Kett00}. Thus a gesture is composed of the physical context in which it is created along with its drawn gestural form. 

More importantly, perhaps, is the distinction between gestures and direct manipulation techniques. Direct manipulation refers to an interface design philosophy which encourages graphical, manual and incremental interactions over more complex, abstract ones. The main premise of this is that it reduces the cognitive demands placed on users by the interface~\citep{Froh93}. \citet{Shne82} provided three properties which can be used to characterise a direct manipulation system. These are a perpetual representation of the objects of interest to the user; labelled buttons or physical actions (for example, drag and drop) instead of complex syntax and ``rapid, incremental reversible operations whose impact on the object of interest is immediately visible"~\citep{Froh93}. This can be further explained as an attempt to reduce the perceived distance between the computerised representation of a subject and the way in which a user might normally think about it. This is often accomplished through the use of real-world metaphors which the user is able to manipulate graphically rather than through a series of complex commands~\citep{Hutc85}. Thus direct manipulation interfaces can be seen to reduce the cognitive load required of users. However, unlike gestural interaction, these interfaces do not provide the ability to rapidly execute complex and abstract commands due to their incremental nature~\citep{Froh93}.

\subsection{Benefits and limitations of gestural interaction}
When considering novel interaction techniques, the benefits and limitations of such interactions should be carefully studied. This is of added importance given the highly abstract nature of gestures.

As discussed above, gestures are often used to encapsulate symbolic meaning. This implies that gestures can be used as a form of shortcut, where multiple commands are specified through the use of a single gesture~\citep{Niel04}. This, in turn, increases communication bandwidth as one is able to complete a number of commands far more quickly. In addition, the context-sensitive nature of gestures means that they can be used to create alternative communication channels within different situations~\citep{Jaim07}.

The main concern regarding the use of gestures is that the purpose or form of gestures may not always be immediately apparent. That is, users may not be easily able to identify what task a gesture performs in relation to its form~\citep{Wobb09}. This is further complicated by the context-sensitivity of gestures. That is, the same gesture may be used for multiple scenarios with its action being determined by the context within which it is performed~\citep{Wobb09}. The reuse of gestures is directly related to a second concern regarding their integration. That is, the amount of learning required of the user. It may not be possible to provide the same affordances, that is, features of the object or interface that suggest how it should be interacted with, in gestural interfaces as are available in direct manipulation designs~\citep{Niel04}. This implies that users will be required to learn and remember a fair number of gestural commands. This, in turn, increases the cognitive demands placed upon the user and may ultimately decrease the usability of the system~\citep{Niel04}.

Thus, while gestures have the potential to provide a means for executing a number of commands rapidly and with a high level of abstraction, care must be taken when selecting gestures for use. Gestures should be easily memorable and provide as many affordances as possible in order to lessen the cognitive requirements placed on the user.

\subsection{Guidelines for gesture selection}
In light of the issues identified above, a number of guidelines have been identified by \citet{Niel04} in order to attempt to make gestural interfaces more usable. These are briefly discussed below.

The first suggestion is that gestures should be memorable and easy to perform~\citep{Niel04}. Related to this, gestures should also be ergonomic, so that unnecessary physical stress is not placed on the user~\citep{Wobb09}. Gestures should also be intuitive and should be used to represent logical metaphorical and iconic links with their proposed functionality. Finally, the system should be able to unambiguously identify all gestures in use~\citep{Jaim07}.

These guidelines primarily serve to decrease the learning and memory demands placed upon users of such interfaces. This in turn, should make gestural interfaces more usable and thus more useful.

\subsection{Conclusion}
Gestural interfaces provide an alternative to the well-established direct manipulation form. However, as is the case for any interface, gestures are not appropriate for all applications. Instead, one needs to consider the specific requirements of the application under development. Thus, while they should not be considered as a replacement for direct manipulation, there are cases where the inclusion of gestural controls can indeed make the manipulation of interface elements far more effective.

